\hypersetup{pageanchor=false}
\pagenumbering{gobble}

\dcover

\dsignatures

\dlibrary{ficha.pdf}
	
\ddedicatory{
	\raggedleft \normalsize Seção reservada à dedicatória.
}

\dacknowledgment{
	Seção reservada a agradecimentos.
	
	Agradece-se à CAPES pelo financiamento parcial do projeto.
}

\dresumo{Esta dissertação analisa a viabilidade técnica, econômica e ambiental do biobunker, definido como a mistura de biodiesel com bunker fuel, como combustível de transição no transporte marítimo. O estudo insere-se no contexto das metas de descarbonização estabelecidas pela Organização Marítima Internacional para 2030 e da necessidade de avaliar alternativas capazes de reduzir emissões no curto e médio prazo, enquanto combustíveis de emissão zero amadurecem tecnologicamente. A análise baseia-se em simulações com dados reais de operação de navios, em estimativas de impactos ambientais derivadas de estudos de Avaliação de Ciclo de Vida (ACV) disponíveis na literatura e na avaliação de custos e restrições logísticas associadas ao uso do biobunker. Adicionalmente, o trabalho examina o papel do arcabouço institucional e das políticas públicas brasileiras na adoção desse combustível pelo setor marítimo. Ao investigar o biobunker como alternativa intermediária na transição energética, a dissertação contribui para a compreensão de suas possibilidades e limitações no contexto do transporte marítimo brasileiro.}{bunker fuel, biodiesel, biobunker, descarbonização, transporte marítimo}

\dabstract{This dissertation analyzes the technical, economic, and environmental feasibility of biobunker, defined as a blend of biodiesel and bunker fuel, as a transitional fuel in maritime transport. The study is situated within the context of the decarbonization targets established by the International Maritime Organization for 2030 and the need to assess alternatives capable of reducing emissions in the short and medium term while zero-emission fuels continue to mature technologically. The analysis is based on simulations conducted using real ship operational data, the use of environmental impact estimates derived from Life Cycle Assessment (LCA) studies available in the literature, and the evaluation of costs and logistical constraints associated with biobunker deployment. In addition, the dissertation examines the role of Brazil's institutional framework and public policies in shaping the adoption of this fuel within the maritime sector. By investigating biobunker as an intermediate option in the energy transition, the dissertation contributes to understanding its potential and limitations in the context of Brazilian maritime transport.}{bunker fuel, biodiesel, biobunker, decarbonization, maritime transport}

\dtables
\hypersetup{pageanchor=true}
